\subsection{Proper retracts across an arbitrary support cut}

\begin{theorem}[Proved: a cut retraction preserves projective-plane type]
\label{CP2:A-cut-retract}
Let $C$ be a finite category whose objects are the disjoint union of
full subcategories $D,E$, with no arrows from $E$ to $D$.
Write $S(d,e)=C(d,e)$ for its connecting profunctor, a functor
$D^{\mathrm{op}}\times E\to\mathbf{FinSet}$.
Suppose $r:S\to S$ is an idempotent natural transformation.
Keep $D,E$ and replace each $S(d,e)$ by $rS(d,e)$, obtaining the
finite subcategory $Y$. If $BC\simeq\CP^2$, then $BY\simeq\CP^2$.
If any component of $r$ is not the identity, $Y$ has strictly fewer
arrows and the same objects and support as $C$.
\end{theorem}
\begin{proof}
Naturality means that every $r_{d,e}$ commutes with precomposition
by arrows of $D$ and postcomposition by arrows of $E$. Its images
therefore form a subprofunctor and define $Y$. There cannot be two
composable crossing arrows. Hence the maps which are identities
on $D,E$ and equal to $r$ on crossing arrows give an actual functor
$R:C\to Y$. For the inclusion $i:Y\hookrightarrow C$, idempotence
gives $Ri=\operatorname{id}_Y$.

Use the natural two-block mapping-cylinder model of
Lemma~\ref{CP2:A-two-block}. Its bridge spaces for $Y$ and $C$ are
$T_Y\subset T_C$, obtained by taking the two-sided bars of the
respective profunctors and then realizing. They are CW subcomplexes.
The two corner spaces are identical. Collapsing the model of $BY$
in that of $BC$ consequently gives
\[
 \frac{T_C\times I}
 {T_C\times\partial I\;\cup\;T_Y\times I}
   =\Sigma(T_C/T_Y).
\]
The models and their inclusions are natural in the profunctor, so
this computes the relative cohomology and its cup product for
$(BC,BY)$. Every product of positive-degree relative classes is
zero, because the displayed quotient is a suspension.

Assume $BC\simeq\CP^2$. The retraction makes each
$H_j(BY;\mathbb Z)$ a direct summand of $H_j(BC;\mathbb Z)$.
Thus $BY$ is connected, its only possible positive integral
homology groups are summands of $\mathbb Z$ in degrees two and
four, and its fundamental group is trivial as a retract of
$\pi_1(BC)=1$. If $H^2(BY;\mathbb Z)=0$, the generator
$u\in H^2(BC;\mathbb Z)$ lifts to relative cohomology; its square
would vanish, a contradiction. Therefore $H^2(BY;\mathbb Z)$ is
$\mathbb Z$ and the split injection
$R^*:H^2(BY;\mathbb Z)\to H^2(BC;\mathbb Z)$ is an isomorphism.
Choose $v$ with $R^*v=u$. Then $R^*(v^2)=u^2$ is the integral
four-dimensional generator. The split degree-four map is also
an isomorphism, and $v^2$ is primitive. The integral recognition
argument of Theorem~\ref{A:monoid-upper23} gives $BY\simeq\CP^2$.

A finite idempotent self-map is a bijection only when it is the
identity. Thus a nonidentity component has a proper image.
Each originally nonempty component remains nonempty, so support
and object set are preserved while the total arrow count drops.
\end{proof}

\begin{corollary}[Proved: endomorphisms of a minimal connecting profunctor]
\label{CP2:A-minimal-cut-endomorphisms}
In a category realizing $\CP^2$ with the least possible total
number of arrows, every natural endomorphism of every connecting
profunctor across a support cut is a componentwise bijection.
The same holds for an acyclic model with a fixed number of objects
and the least arrow count among such models.
\end{corollary}
\begin{proof}
For a natural endomorphism $f$, choose a common positive exponent
$k$ which is at least the transient length of every finite component
map and is divisible by every eventual cycle length. Then $f^k$
is a natural idempotent. If any component of $f$ is not bijective,
its eventual image is proper and that component of $f^k$ differs
from the identity. The theorem gives a strictly smaller category
of the same homotopy type. In the acyclic case, the object set,
trivial corners and support are unchanged, so it is also a smaller
acyclic model at the same object count. Both contradict minimality.
\end{proof}
